% ═══════════════════════════════════════════════════════════════
% LESEBLATT - Physik Klasse 6: Bewegungsarten (Referenz)
% Kompakte Übersicht zum Nachschlagen
% ═══════════════════════════════════════════════════════════════

\documentclass[11pt, a4paper]{article}

\usepackage[utf8]{inputenc}
\usepackage[T1]{fontenc}
\usepackage[ngerman]{babel}
\usepackage{helvet}
\renewcommand{\familydefault}{\sfdefault}

\usepackage[margin=1.5cm, top=1.8cm, bottom=1.5cm]{geometry}
\usepackage{enumitem}
\usepackage{tikz}
\usepackage{fancyhdr}
\usepackage{xcolor}
\usepackage{tcolorbox}
\tcbuselibrary{skins}

% Farben
\definecolor{physik}{HTML}{2c5aa0}
\definecolor{physik-hell}{HTML}{e8f0fa}

% Kopfzeile
\pagestyle{fancy}
\fancyhf{}
\fancyhead[L]{\small\textcolor{physik}{\textbf{Physik}} | Bewegungsarten -- Referenz}
\fancyhead[R]{\small Klasse 6}
\renewcommand{\headrulewidth}{0pt}

% Merkkasten
\newtcolorbox{merkkasten}[1][]{
    colback=physik-hell,
    colframe=physik,
    boxrule=1.5pt,
    arc=0pt,
    left=8pt, right=8pt, top=6pt, bottom=6pt,
    #1
}

\setlength{\parindent}{0pt}
\setlength{\parskip}{4pt}

\begin{document}

% ═══════════════════════════════════════════════════════════════
% TITEL
% ═══════════════════════════════════════════════════════════════

\begin{center}
{\Large\bfseries\textcolor{physik}{Bewegungsarten -- Übersicht}}
\end{center}
\vspace{0.3em}
\hrule
\vspace{0.8em}

% ═══════════════════════════════════════════════════════════════
% DEFINITION
% ═══════════════════════════════════════════════════════════════

\begin{merkkasten}[title={\small\bfseries\textcolor{white}{Definition}}, coltitle=white, attach boxed title to top left={yshift=-2mm, xshift=5mm}, boxed title style={colback=physik, arc=0pt}]
Ein Körper ist in \textbf{Bewegung}, wenn er seinen \textbf{Ort} verändert.

Wir beschreiben Bewegungen mit:
\begin{itemize}[leftmargin=1.5em, itemsep=1pt, topsep=2pt]
    \item \textbf{Bahnform} -- Wie sieht der Weg aus?
    \item \textbf{Bewegungsart} -- Ändert sich die Geschwindigkeit?
\end{itemize}
\end{merkkasten}

\vspace{0.5em}

% ═══════════════════════════════════════════════════════════════
% ZWEI SPALTEN
% ═══════════════════════════════════════════════════════════════

\begin{minipage}[t]{0.48\textwidth}

\begin{merkkasten}[title={\small\bfseries\textcolor{white}{Bahnformen}}, coltitle=white, attach boxed title to top left={yshift=-2mm, xshift=5mm}, boxed title style={colback=physik, arc=0pt}]

\textbf{Geradlinig}\\
Der Körper bewegt sich auf einer \textbf{geraden Linie}.

\vspace{0.3em}
\begin{center}
\begin{tikzpicture}[scale=0.7]
    \draw[thick, physik, ->] (0,0) -- (3,0);
    \node[below, font=\scriptsize] at (1.5,-0.2) {Auto auf Autobahn};
\end{tikzpicture}
\end{center}

\vspace{0.3em}
\textbf{Krummlinig}\\
Der Körper bewegt sich auf einer \textbf{gebogenen Bahn}.

\vspace{0.3em}
\begin{center}
\begin{tikzpicture}[scale=0.7]
    \draw[thick, physik, ->] (0,0) arc (180:-90:0.6);
    \node[below, font=\scriptsize] at (0.6,-0.8) {Karussell};
    \draw[thick, physik, ->] (2.5,0) .. controls (3,0.5) and (3.5,-0.3) .. (4,0.2);
    \node[below, font=\scriptsize] at (3.25,-0.3) {Schlange};
\end{tikzpicture}
\end{center}

\vspace{0.2em}
{\small\textit{Kreisförmig = Sonderfall von krummlinig}}

\end{merkkasten}

\end{minipage}
\hfill
\begin{minipage}[t]{0.48\textwidth}

\begin{merkkasten}[title={\small\bfseries\textcolor{white}{Bewegungsarten}}, coltitle=white, attach boxed title to top left={yshift=-2mm, xshift=5mm}, boxed title style={colback=physik, arc=0pt}]

\textbf{Gleichförmig}\\
Die Geschwindigkeit bleibt \textbf{gleich}.

\vspace{0.3em}
\begin{center}
\begin{tikzpicture}[scale=0.7]
    \foreach \x in {0,0.8,1.6,2.4} {
        \fill[physik] (\x,0) circle (3pt);
    }
    \draw[->] (0,-0.4) -- (2.8,-0.4);
    \node[below, font=\scriptsize] at (1.2,-0.5) {gleiche Abstände};
\end{tikzpicture}
\end{center}

\vspace{0.3em}
\textbf{Ungleichförmig}\\
Die Geschwindigkeit \textbf{ändert} sich.

\vspace{0.3em}
\begin{center}
\begin{tikzpicture}[scale=0.7]
    \foreach \x in {0,0.2,0.6,1.2,2.0} {
        \fill[physik] (\x,0) circle (3pt);
    }
    \draw[->] (0,-0.4) -- (2.4,-0.4);
    \node[below, font=\scriptsize] at (1,-0.5) {verschiedene Abstände};
\end{tikzpicture}
\end{center}

\vspace{0.2em}
{\small\textit{z.B. Anfahren, Bremsen}}

\end{merkkasten}

\end{minipage}

\vspace{0.8em}

% ═══════════════════════════════════════════════════════════════
% BEISPIELE
% ═══════════════════════════════════════════════════════════════

\begin{merkkasten}[title={\small\bfseries\textcolor{white}{Beispiele}}, coltitle=white, attach boxed title to top left={yshift=-2mm, xshift=5mm}, boxed title style={colback=physik, arc=0pt}]
\begin{center}
\renewcommand{\arraystretch}{1.3}
\begin{tabular}{|l|c|c|}
\hline
\textbf{Beispiel} & \textbf{Bahnform} & \textbf{Bewegungsart} \\
\hline
Zug auf gerader Strecke & geradlinig & gleichförmig \\
\hline
Fahrrad beim Anfahren & geradlinig & ungleichförmig \\
\hline
Gondel am Riesenrad & krummlinig (kreisf.) & gleichförmig \\
\hline
Achterbahn in der Kurve & krummlinig & ungleichförmig \\
\hline
Minutenzeiger einer Uhr & krummlinig (kreisf.) & gleichförmig \\
\hline
\end{tabular}
\end{center}
\end{merkkasten}

\end{document}
